%!TEX root=./neurips_2023.tex

\section{Discussion}
% \jh{limitation: arch needs to be the same, compared to adapterfusion}
This study aims to compose trained parameter-efficient modules (PEMs) in parameter space, utilizing linear arithmetic, to create a highly adaptable manipulation of the module capabilities.
We introduce addition and negation operators for the PEM serving as the fundamental operators. 
We combine them to execute flexible linear arithmetic operations on the module parameters to attain various objectives. These objectives involve aggregating PEMs together for distribution generalization and to facilitate multi-tasking, negating for unlearning certain skills, and combining PEMs of related domains and tasks for domain transfer.
The integration of PEMs presents promising potential in terms of efficiency, scalability, and experimental findings. 
Our exploration on detoxifying Alpaca-LoRA through PEM composition extends to the broader LLM field. 
% The larger scale of LLM and higher training costs of full fine-tuning leads to affordable parameter-efficient tuning more abundant, this study is becoming increasingly necessary and appealing.
\paragraph{Potential Impacts and Limitations:}
Our work on composing existing PEMs may inherit the biases or safety concerns that inherently exist in these PEMs. 
Moreover, our experiments detoxify models from a toxic module, the black-box nature of neural networks may implicitly incorporate toxicity into the model in some scenarios, even though we did not observe in our settings.
Limitations of this work include (1) we restricted the exploration to the identical PEM architecture, and the same module initialization in most of the experiments; and (2) our approach requires tuning the weight hyperparameter $\lambda$. Future work will focus on exploring alternative composition of PEMs with different architectures and varied module initialization, and computing the weight hyperparamter through automatic methods as in~\citet{jin2022dataless}.
% which holds immense potential for the integration of PEMs in diverse fields.